% !TEX program = xelatex
\documentclass[11pt,a4paper]{article}

\usepackage{fontspec}
\usepackage{polyglossia}
\usepackage{amsmath,amssymb,amsfonts}
\usepackage{xcolor}
\usepackage{graphicx}
\usepackage{geometry}
\usepackage{tikz}
\usetikzlibrary{positioning,arrows.meta,shapes}

\geometry{left=2cm,right=2cm,top=2cm,bottom=2cm}

\setdefaultlanguage{persian}
\setotherlanguage{english}
\setmainfont{XB Niloofar}

\newcommand{\en}[1]{\textenglish{#1}}
\newcommand{\separator}{\vspace{1em}\noindent\rule{\linewidth}{0.4pt}\vspace{1em}}
\newcommand{\thickseparator}{\vspace{1em}\noindent\rule{\linewidth}{1pt}\vspace{1em}}

\title{\textbf{پاسخ سوال ۱} \\[0.4cm] \small معماری شبکه عصبی و الگوریتم \en{Backpropagation}}
\author{محمد باغبان زاده \\ شماره دانشجویی: شناسهٔ دانشجو}
\date{}
\linespread{1.35}

\begin{document}

\maketitle
\thickseparator

\section{مفاهیم پایه: معماری شبکه عصبی (فصل \en{۶} کتاب)}

بر اساس بخش‌های \textbf{\en{۶-۱}}، \textbf{\en{۶-۳}} و \textbf{\en{۶-۴}} کتاب، شبکه عصبی ارائه شده در سوال دارای معماری زیر است:

\begin{itemize}
    \item \textbf{لایه ورودی (\en{Input Layer})}: یک نورون \(x_1\)
    \item \textbf{لایه پنهان (\en{Hidden Layer})}: دو نورون \(z_1\) و \(z_2\) با تابع فعال‌سازی \textbf{\en{ReLU}} (بخش \en{۶-۱}، معادله \en{۶-۶})
    \item \textbf{لایه خروجی (\en{Output Layer})}: یک نورون با تابع فعال‌سازی \textbf{\en{Sigmoid}} (بخش \en{۶-۱}، معادله \en{۶-۳})
\end{itemize}

طبق شکل سوال، وزن‌ها و بایاس‌ها به صورت زیر هستند:

\begin{align*}
&\text{:وزن‌های لایه پنهان} \quad w_{11} = 0.2,\; w_{12} = -0.1,\; w_{21} = 0.7,\; w_{22} = 0.7 \\
&\text{:بایاس‌های لایه پنهان} \quad b_{11} = 0.04,\; b_{12} = 0.01 \\
&\text{:بایاس لایه خروجی} \quad b_{21} = 0.08
\end{align*}

\separator
\section{قسمت (الف): پیش‌بینی کلاس برای \(x_1 = 0.3\)}

\subsection{محاسبه \en{Forward Pass}}

ابتدا خروجی لایه پنهان را محاسبه می‌کنیم. برای هر نورون، ابتدا مجموع وزنی را محاسبه کرده سپس تابع \en{ReLU} را اعمال می‌کنیم:

\[
\text{ReLU}(z) = \max(0, z)
\]

\textbf{نورون \(z_1\):}
\[
z_{1,\text{net}} = w_{11} \cdot x_1 + b_{11} = 0.2 \times 0.3 + 0.04 = 0.06 + 0.04 = 0.10
\]
\[
z_1 = \text{ReLU}(0.10) = 0.10
\]

\textbf{نورون \(z_2\):}
\[
z_{2,\text{net}} = w_{12} \cdot x_1 + b_{12} = (-0.1) \times 0.3 + 0.01 = -0.03 + 0.01 = -0.02
\]
\[
z_2 = \text{ReLU}(-0.02) = 0
\]

سپس خروجی لایه نهایی (لایه خروجی) را با تابع سیگموئید محاسبه می‌کنیم. طبق شکل، وزن‌های اتصال از لایه پنهان به خروجی \(w_{21} = 0.7\) و \(w_{22} = 0.7\) و بایاس خروجی \(b_{21} = 0.08\) است:

\[
y_{\text{net}} = w_{21} \cdot z_1 + w_{22} \cdot z_2 + b_{21} = 0.7 \times 0.10 + 0.7 \times 0 + 0.08 = 0.07 + 0 + 0.08 = 0.15
\]

\[
\hat{y} = \sigma(y_{\text{net}}) = \frac{1}{1 + e^{-0.15}}
\]

مقدار \(e^{-0.15} \approx 0.8607\)، بنابراین:

\[
\hat{y} \approx \frac{1}{1 + 0.8607} = \frac{1}{1.8607} \approx 0.5374
\]

\subsection{قاعده تصمیم‌گیری}

با آستانه ۰.۵، اگر \(\hat{y} \ge 0.5\) کلاس ۱ و در غیر این صورت کلاس ۰ پیش‌بینی می‌شود. از آنجا که \(0.5374 > 0.5\):

\[
\boxed{\text{کلاس پیش‌بینی‌شده: } 1}
\]

\separator
\section{قسمت (ب): محاسبه \en{Backpropagation} برای \(w_{12}\) و \(b_{12}\)}

برای محاسبه گرادیان‌ها، فرض می‌کنیم تابع هزینه از نوع \textbf{باینری کراس-آنتروپی} است (بخش \en{۶-۶-۱} کتاب، معادله \en{۶-۲۵}):

\[
L(\hat{y}, y) = -[y \log \hat{y} + (1-y) \log(1-\hat{y})]
\]

\subsection{گام ۱: محاسبه سیگنال خطا در لایه خروجی}

مشتق تابع هزینه نسبت به خروجی شبکه (\(\hat{y}\)) برای یک نمونه با برچسب واقعی \(y\) به صورت زیر است:

\[
\delta^{(3)} = \frac{\partial L}{\partial y_{\text{net}}} = \hat{y} - y
\]

(این فرمول از مشتق‌گیری تابع کراس-آنتروپی و سیگموئید به دست می‌آید، همان‌طور که در بخش \en{۶-۶-۴} کتاب، معادله \en{۶-۴۳} نشان داده شده است.)

در این سوال، برچسب واقعی \(y\) داده نشده است. معمولاً در مسائل طبقه‌بندی، فرض می‌شود نمونه متعلق به کلاس مثبت (\(y=1\)) است. با این فرض:

\[
\delta^{(3)} = \hat{y} - y = 0.5374 - 1 = -0.4626
\]

\subsection{گام ۲: انتشار خطا به لایه پنهان}

سیگنال خطا برای نورون \(z_2\) در لایه پنهان به صورت زیر محاسبه می‌شود:

\[
\delta^{(2)}_2 = \delta^{(3)} \cdot w_{22} \cdot \text{ReLU}'(z_{2,\text{net}})
\]

که در آن \(\text{ReLU}'(z)\) مشتق تابع \en{ReLU} است (بخش \en{۶-۶-۴} کتاب، معادله \en{۶-۴۰}):

\[
\text{ReLU}'(z) = 
\begin{cases} 
1 & \text{\en{if }} z > 0 \\
0 & \text{\en{if }} z \le 0
\end{cases}
\]

از آنجا که \(z_{2,\text{net}} = -0.02 \le 0\)، داریم \(\text{ReLU}'(-0.02) = 0\). بنابراین:

\[
\delta^{(2)}_2 = (-0.4626) \times 0.7 \times 0 = 0
\]

\subsection{گام ۳: محاسبه گرادیان نسبت به \(w_{12}\) و \(b_{12}\)}

گرادیان تابع هزینه نسبت به وزن \(w_{12}\) و بایاس \(b_{12}\) به صورت زیر محاسبه می‌شود:

\[
\frac{\partial L}{\partial w_{12}} = \delta^{(2)}_2 \cdot x_1 = 0 \times 0.3 = 0
\]

\[
\frac{\partial L}{\partial b_{12}} = \delta^{(2)}_2 \cdot 1 = 0 \times 1 = 0
\]

\subsection{نتیجه}

از آنجا که \(z_{2,\text{net}} < 0\)، نورون \(z_2\) در ناحیه صفر تابع \en{ReLU} قرار دارد و مشتق آن صفر است. بنابراین سیگنال خطا در این نورون خنثی شده و به نورون‌های قبلی منتقل نمی‌شود. در نتیجه:

\[
\boxed{\frac{\partial L}{\partial w_{12}} = 0 \quad \text{و} \quad \frac{\partial L}{\partial b_{12}} = 0}
\]

\separator
\section{جمع‌بندی و پاسخ نهایی}

\begin{center}
\fbox{\parbox{0.9\linewidth}{\centering
\textbf{پاسخ نهایی:}
\begin{itemize}
    \item \textbf{قسمت (الف):} با \(x_1 = 0.3\)، خروجی مدل \(\hat{y} \approx 0.5374 > 0.5\) است. بنابراین شبکه \textbf{کلاس ۱} را پیش‌بینی می‌کند.
    \item \textbf{قسمت (ب):} 
    \[
    \frac{\partial L}{\partial w_{12}} = 0, \qquad \frac{\partial L}{\partial b_{12}} = 0
    \]
    دلیل: چون \(z_{2,\text{net}} = -0.02 < 0\)، مشتق \en{ReLU} صفر است و گرادیان به این وزن و بایاس نمی‌رسد.
\end{itemize}
}}
\end{center}

\thickseparator

\section{نکات تکمیلی}

\begin{itemize}
    \item بایاس لایه خروجی \(b_{21} = 0.08\) در محاسبات لحاظ شده است.
    \item اگر برچسب واقعی نمونه \(y=0\) بود، \(\delta^{(3)} = 0.5374 - 0 = 0.5374\) می‌شد که همچنان به دلیل صفر بودن مشتق \en{ReLU}، گرادیان‌ها صفر باقی می‌ماند.
    \item این مثال نشان می‌دهد که چگونه نورون‌های خاموش (\en{dead neurons}) در شبکه‌های عمیق می‌توانند باعث قطع شدن جریان گرادیان شوند (بخش \en{۶-۶} کتاب).
\end{itemize}

\thickseparator

\end{document}